\chapter{Plasmons, Polaritons, and Polarons}

This chapter treats three types of collective excitations that involve the coupling of electromagnetic fields with charged particles in solids: plasmons (oscillations of the electron gas), polaritons (coupling of photons with optical phonons), and polarons (electrons dressed by a cloud of phonons).

\section{Dielectric Function of the Electron Gas}

The free electron gas has a frequency-dependent dielectric function:
\begin{equation}
  \varepsilon(\omega) = 1 - \frac{\omega_p^2}{\omega^2}
\end{equation}
where the \textbf{plasma frequency} is:
\begin{equation}
  \omega_p = \sqrt{\frac{ne^2}{m\varepsilon_0}}
\end{equation}

\begin{table}[H]
\centering
\caption{Plasma frequencies and corresponding wavelengths for free-electron metals.}
\label{tab:plasma}
\begin{tabular}{@{}lSSS@{}}
\toprule
Metal & {$n$ ($10^{28}$ \si{\per\metre\cubed})} & {$\hbar\omega_p$ (\si{\electronvolt})} & {$\lambda_p$ (\si{\nano\metre})} \\
\midrule
Li & 4.70 & 8.0  & 155 \\
Na & 2.65 & 5.9  & 210 \\
K  & 1.40 & 4.3  & 290 \\
Al & 18.1 & 15.8 & 79 \\
Cu & 8.47 & 10.8 & 115 \\
Ag & 5.86 & 9.0  & 138 \\
\bottomrule
\end{tabular}
\end{table}

\subsection{UV Transparency of Metals}

For $\omega > \omega_p$, $\varepsilon > 0$ and electromagnetic waves propagate---the metal becomes transparent. The alkali metals are transparent in the ultraviolet, with measured transparency edges in good agreement with the free-electron $\omega_p$. Sodium becomes transparent at $\lambda \approx \SI{210}{\nano\metre}$, potassium at $\lambda \approx \SI{315}{\nano\metre}$.

\subsection{Plasmons: Electron Energy Loss Spectroscopy}

Plasmons are quantized plasma oscillations with energy $\hbar\omega_p$. They are directly observed in \textbf{electron energy loss spectroscopy (EELS)}: a beam of fast electrons passing through a thin metal film loses energy in discrete multiples of $\hbar\omega_p$.

Wu et al.\ (2017) reviewed EELS measurements, noting that the bulk plasmon of aluminum at $\hbar\omega_p \approx \SI{15}{\electronvolt}$ was first observed by Ruthemann in 1948. Modern EELS with sub-\si{\electronvolt} resolution can map plasmon modes with nanometer spatial resolution.


\section{Polaritons}

When infrared light propagates through an ionic crystal near a transverse optical (TO) phonon frequency, the photon and phonon couple to form a mixed mode called a \textbf{polariton}. The dielectric function near the TO phonon is:
\begin{equation}
  \varepsilon(\omega) = \varepsilon_\infty \frac{\omega^2 - \omega_{LO}^2}{\omega^2 - \omega_{TO}^2}
\end{equation}

This gives a frequency range $\omega_{TO} < \omega < \omega_{LO}$ where $\varepsilon < 0$ and light cannot propagate---the \textbf{Reststrahlen band}. The ratio of LO to TO frequencies is given by the \textbf{Lyddane--Sachs--Teller relation}:
\begin{equation}
  \frac{\omega_{LO}^2}{\omega_{TO}^2} = \frac{\varepsilon_0}{\varepsilon_\infty}
\end{equation}

\begin{table}[H]
\centering
\caption{TO and LO phonon frequencies and dielectric constants for ionic crystals, verifying the LST relation.}
\label{tab:lst}
\begin{tabular}{@{}lSSSSS@{}}
\toprule
Crystal & {$\omega_{TO}$ (\si{\tera\hertz})} & {$\omega_{LO}$ (\si{\tera\hertz})} & {$\varepsilon_0$} & {$\varepsilon_\infty$} & {$\varepsilon_0/\varepsilon_\infty$ (calc)} \\
\midrule
NaCl  & 4.92 & 7.91 & 5.90 & 2.25 & 2.59 \\
KBr   & 3.39 & 5.07 & 4.90 & 2.33 & 2.24 \\
GaAs  & 8.06 & 8.76 & 12.9 & 10.9 & 1.18 \\
MgO   & 12.1 & 21.5 & 9.83 & 2.95 & 3.16 \\
\bottomrule
\end{tabular}
\end{table}


\section{Polarons}

When an electron moves through an ionic crystal, it polarizes the surrounding lattice, creating a potential well that follows the electron. The electron plus its polarization cloud is called a \textbf{polaron}. The polaron has an enhanced effective mass $m^*_{\text{pol}} > m^*_{\text{band}}$ and a reduced mobility.

The electron--phonon coupling strength is characterized by the Fr\"ohlich coupling constant $\alpha$. For weak coupling ($\alpha < 1$), the mass enhancement is $m^*_{\text{pol}}/m^* \approx 1 + \alpha/6$. For strong coupling ($\alpha > 6$), the polaron becomes self-trapped.

\vspace{1cm}
\noindent\rule{\textwidth}{0.4pt}
\textit{The optical processes of interband transitions, excitons, and Raman scattering are presented in the next chapter.}
